\subsection{Profil Web-cloud}
\label{subsec:profil-web-cloud}

\texttt{web-cloud} deklariert \texttt{frontend}, \texttt{backend} und
\texttt{cloud}. Der Scaffold enthält damit Vite-Frontend, den für das
Basisbackend ausgewählten FastAPI-Code sowie \path{deployment/} und
\path{.dockerignore}; \path{src-tauri/} und dessen npm-Abhängigkeiten werden
nicht übernommen. Das Frontend kommuniziert über die öffentliche API-URL mit
dem getrennten Backend. \texttt{cloud} kennzeichnet Deployment-Dateien ohne
Bindung an einen konkreten Anbieter
\parencite{kleiveist2026deployment,kleiveist2026profiles}.

Ohne weitere Auswahl fehlen Datenbankmodul, Alembic sowie die Datenbank- und
Psycopg-Dateien. PostgreSQL ist deshalb eine kompatible, aber
separat anzufordernde Capability; \texttt{web-cloud} ist nicht gleichbedeutend
mit \texttt{web-cloud + postgres}.
